% -----------------------------------------------
% ISMIR 2026 LBD submission
% -----------------------------------------------

\documentclass{article}
\usepackage[T1]{fontenc}
\usepackage[utf8]{inputenc}
\usepackage[lbd]{ismir}
\usepackage{amsmath,cite,url}
\usepackage{graphicx}
\usepackage{color}
\usepackage{booktabs}
\usepackage{tikz}
\usetikzlibrary{shapes.geometric,arrows.meta,positioning,fit}
\usepackage{amssymb} % \checkmark
\usepackage{pifont} % \ding{55}
\usepackage[hidelinks,breaklinks]{hyperref}
\def\lbd{}

\title{Classifying Irish Tune Types from ABC Notation with a Character-Level LSTM}

\oneauthor
  {Mario Vallejo Reyes}
  {Department of Media and Digital Culture, Tecnológico de Monterrey, Mexico \\ mariovallejo@tec.mx}

\def\authorname{M. Vallejo Reyes}

\sloppy

\begin{document}

\maketitle

\begin{abstract}
Irish traditional tune types (reels, jigs, hornpipes, etc.) have been classified using hand-crafted symbolic features or audio embeddings of synthesized scores; neural sequence models trained on the same ABC-notation corpora have so far been used only for tune generation, not classification. This work trains a character-level LSTM directly on raw ABC-notation transcriptions from The Session to classify tune type across twelve categories, and compares it against a support-vector-machine baseline built on music21-derived pitch, interval, and rhythm histograms, following prior symbolic-feature classifiers for this task. Both models are evaluated on a piece-level train/validation/test split to avoid leakage between settings of the same tune. Without class-weighted training, the LSTM reaches 83.9\% test accuracy, exceeding the SVM baseline's 82.7\%, but the SVM retains a higher macro-F1 (0.739 vs.\ 0.692), reflecting a trade-off between overall accuracy and balanced performance on metrically similar minority classes. A lightweight neural sequence model can match a hand-engineered baseline directly from text, without audio synthesis or manual feature design, while class imbalance remains an open challenge.
\end{abstract}

\section{Introduction}\label{sec:introduction}

Irish traditional tune types are distinguished primarily by meter and rhythmic feel rather than pitch content, and correctly labelling a tune with its type is a routine task for archives and community repositories such as The Session~\cite{TheSession}, which hosts tens of thousands of tunes in ABC notation, a compact text-based symbolic format. Prior work has approached this task with classical classifiers on hand-crafted symbolic features: Martins and Silla trained decision trees reaching 92.3\% overall precision across eleven genres, with substantially lower precision on metrically similar types such as barndance~\cite{Martins2015}, while Kermit-Canfield and Roman found that reels and hornpipes, which share a similar meter, remained close to chance even for the best classifier~\cite{Kermit2015}. Other work has used deep embeddings of audio synthesized from scores~\cite{Jimenez2024}, and, most recently, data augmentation applied directly to symbolic ABC scores~\cite{Navarro2025}. Separately, neural sequence models have been trained on the same ABC corpora, but only for tune generation~\cite{Sturm2015,Wu2023}. To the best of the author's knowledge, no prior work trains a neural sequence model end-to-end and directly on raw ABC text for tune-type classification, without synthesizing audio or engineering features by hand. This work addresses that gap and reports where a character-level LSTM succeeds and fails relative to a feature-based SVM baseline. Code and additional results are available online\footnote{\url{https://mariovallejotec.github.io/irish-tune-classification/}}.

\section{Method}\label{sec:method}

\begin{figure}[t]
\centering
\begin{tikzpicture}[
    font=\scriptsize,
    >=Latex,
    src/.style={draw, rounded corners, align=center, minimum height=0.6cm,
                text width=7.2cm, inner sep=3pt, fill=gray!12},
    path/.style={draw, align=center, minimum height=0.9cm,
                 text width=3.45cm, inner sep=3pt, fill=green!8},
    out/.style={draw, rounded corners, align=center, minimum height=0.6cm,
                text width=7.2cm, inner sep=3pt, fill=gray!12},
    arr/.style={-{Latex[length=1.8mm]}, line width=0.5pt}
  ]
  \node[src] (data) at (0,0) {ABC-notation tune (raw text), from The Session};
  \node[path] (feat) at (-1.85,-1.15) {music21: 60 pitch/rhythm features};
  \node[path] (tok) at (1.85,-1.15) {Character tokenizer (raw text)};
  \node[path] (svm) at (-1.85,-2.3) {SVM (RBF kernel)};
  \node[path] (lstm) at (1.85,-2.3) {BiLSTM (hidden 128)};
  \node[out] (pred) at (0,-3.45) {Tune-type prediction (12 classes)};

  \draw[arr] (data.south) -- ++(0,-0.15) -| (feat.north);
  \draw[arr] (data.south) -- ++(0,-0.15) -| (tok.north);
  \draw[arr] (feat) -- (svm);
  \draw[arr] (tok) -- (lstm);
  \draw[arr] (svm.south) -- ++(0,-0.2) -| (pred.north);
  \draw[arr] (lstm.south) -- ++(0,-0.2) -| (pred.north);
\end{tikzpicture}
\caption{The two classification pipelines compared in this work: a feature-engineered SVM baseline and an end-to-end LSTM trained directly on raw ABC text.}
\label{fig:pipeline}
\end{figure}

Figure~\ref{fig:pipeline} summarizes the two pipelines compared in this work. Because reels, jigs, and hornpipes are set apart mainly by meter and rhythmic feel rather than melodic content, a classifier for this task must learn rhythmic structure directly from the transcription.

\textbf{Data.} Tunes and their type labels are drawn from The Session~\cite{TheSession} (\texttt{tunes.csv}), yielding 55,352 settings across 23,300 unique tunes and 12 tune types after parsing, with a substantial class imbalance (reel: 20,187; three-two: 428). Data are split 70/15/15 into train/validation/test at the tune level, not the setting level, so that no tune appears in more than one split, following the leakage-avoidance practice established for symbolic music classification~\cite{Zhang2023}.

\textbf{Baseline.} Following~\cite{Kermit2015,Martins2015}, a support-vector machine is trained on 60 hand-crafted features per tune extracted with music21~\cite{Cuthbert2010}: pitch-class and melodic-interval histograms, melodic direction statistics, note-duration histograms, and meter.

\textbf{Neural model.} A single-layer bidirectional LSTM (hidden size 128), implemented in PyTorch, is trained directly on the raw ABC character sequence (no feature extraction, max length 600), with a linear classification head. The full model totals 209,804 parameters. Training uses AdamW~\cite{Loshchilov2019} with warmup and cosine decay, gradient clipping, dropout (0.3) and weight decay ($10^{-4}$), and early stopping on validation loss (patience 10; Fig.~\ref{fig:curve}) -- hyperparameter choices informed by common practice across recent neural network papers~\cite{Li2026,Gajecki2026,Zohar2026}, since no prior work addresses class imbalance for a neural sequence model on this exact task.

\begin{figure}[t]
\centering
\includegraphics[width=\columnwidth]{training_curve.png}
\caption{LSTM train/validation loss by epoch. Validation loss reaches its minimum at epoch 10, after which training loss keeps decreasing while validation loss rises.}
\label{fig:curve}
\end{figure}

\section{Results}\label{sec:results}

\begin{table}[h]
\begin{center}
\begin{tabular}{lcc}
\toprule
Model & Accuracy & Macro-F1 \\
\midrule
SVM (RBF) & 82.7\% & \textbf{0.739} \\
LSTM (class-weighted) & 78.4\% & 0.653 \\
LSTM (unweighted) & \textbf{83.9\%} & 0.692 \\
\bottomrule
\end{tabular}
\end{center}
\caption{Test-set results for the SVM baseline and the LSTM under class-weighted and unweighted training.}
\label{tab:results}
\end{table}

\begin{figure}[t]
\centering
\includegraphics[width=\columnwidth]{confusion_matrix.png}
\caption{Confusion matrix of the unweighted LSTM on the test set (row-normalized).}
\label{fig:confmat}
\end{figure}

The unweighted LSTM's higher accuracy comes from strong performance on majority classes (reel F1 0.882, jig F1 0.966), while precision and recall collapse on tune types that share a meter with more frequent classes (barndance: precision 0.383, recall 0.303; march: precision 0.505, recall 0.302). This is not an arbitrary error pattern (Fig.~\ref{fig:confmat}): reel, hornpipe, and barndance are all played in 4/4 and share a similar eighth-note pulse, whereas jig's distinct 6/8 compound feel gives the model an unambiguous rhythmic cue -- the same reel/hornpipe-type confusions reported for the SVM baseline in~\cite{Kermit2015}. Class-weighted training raises minority-class recall but lowers overall accuracy, indicating that class imbalance, not architecture choice, is the main remaining obstacle for this task.

\textbf{Ablation.} Table~\ref{tab:ablation} evaluates the contribution of class-weighted loss, regularization (dropout and weight decay), and symbolic transposition augmentation~\cite{Navarro2025} to the LSTM's performance. Augmentation lowers both accuracy and macro-F1 relative to the unaugmented model. Regularization alone leaves test performance almost unchanged while class weighting is on; only removing class weighting, on top of regularization, produces a substantial shift toward higher accuracy at the cost of macro-F1.

\begin{table}[h]
\begin{center}
\begin{tabular}{ccc|cc}
\toprule
Class w. & Reg. & Aug. & Accuracy & Macro-F1 \\
\midrule
\checkmark & \ding{55} & \ding{55} & 77.5\% & 0.650 \\
\checkmark & \ding{55} & \checkmark & 76.5\% & 0.628 \\
\checkmark & \checkmark & \ding{55} & 78.4\% & 0.653 \\
\ding{55} & \checkmark & \ding{55} & \textbf{83.9\%} & \textbf{0.692} \\
\bottomrule
\end{tabular}
\end{center}
\caption{Ablation over the LSTM training recipe: class-weighted loss (Class w.), regularization (Reg.: dropout + weight decay), and transposition augmentation (Aug.).}
\label{tab:ablation}
\end{table}

\section{Conclusion}\label{sec:conclusion}

A character-level LSTM trained directly on raw ABC-notation text can match or exceed a feature-engineered SVM baseline in accuracy on Irish tune-type classification, without audio synthesis or hand-crafted features, though it does not resolve the class-imbalance problem that the same baseline family already exhibited. Future work includes cross-validation, comparison against recent symbolic-augmentation approaches~\cite{Navarro2025}, and dedicated class-imbalance techniques for sequence models.

\bibliography{ISMIRtemplate}

\end{document}
